%! Confidence Intervals for the Difference of Two Means — Unit 7, Lesson 7.6 — Group Activity
\documentclass[10pt]{article}
\usepackage{apstats-article}
\usepackage{apstats-boxes}
\newcommand{\TallMath}[1]{$\displaystyle #1\rule[-1.4em]{0pt}{3.2em}$}

\begin{document}

\pageheader{Unit 7, Lesson 7.6}{Group Activity}
\namepartnerperiod

\vspace{0.15in}

\textit{For each scenario: verify conditions, compute the standard error, find the confidence
interval, and interpret the result. Show all work.}

\begin{scenariobox}[Scenario 1 --- Sleep Hours: Teens vs.\ Adults (Tier R)]{sky}
A health researcher randomly selects 32 teenagers and 35 adults and records their average
nightly sleep (hours). Results:
Teens: $\bar x_1 = 7.1$ hr, $s_1 = 1.2$ hr. \quad
Adults: $\bar x_2 = 7.8$ hr, $s_2 = 0.9$ hr.
Both sample distributions are roughly symmetric with no outliers.

\begin{enumerate}[label=\alph*.]
  \item \textbf{Identify the procedure} and write the point estimate.

        Procedure: \blank{8cm}

        Point estimate: $\bar x_1 - \bar x_2 = $ \blank{2cm} $-$ \blank{2cm} $=$ \blank{2.5cm}

  \item \textbf{Verify conditions.}
        \begin{itemize}
          \item Random: \writeline
          \item Normal: \writeline
        \end{itemize}

  \item \textbf{Compute} the standard error. Use the template:

        $SE = \sqrt{\dfrac{s_1^2}{n_1} + \dfrac{s_2^2}{n_2}}
            = \sqrt{\dfrac{(\blank{1.5cm})^2}{\blank{1.5cm}} + \dfrac{(\blank{1.5cm})^2}{\blank{1.5cm}}}
            = \sqrt{\blank{3cm} + \blank{3cm}} \approx$ \blank{2cm}

  \item Using technology: $df \approx 61$, $t^* \approx 2.000$ (95\% confidence).

        $ME = t^* \times SE = 2.000 \times \blank{2cm} \approx$ \blank{2cm}

        $CI = ($ \blank{2cm} $\pm$ \blank{2cm} $) = ($ \blank{2.5cm} , \blank{2.5cm} $)$

  \item \textbf{Interpret} the interval in context.
        \writelines{2}

  \item Does the interval contain 0? \blank{1.5cm} \quad
        What does this suggest about $\mu_{\text{teens}} - \mu_{\text{adults}}$? \writeline
\end{enumerate}
\end{scenariobox}

\begin{scenariobox}[Scenario 2 --- Screen Time: Urban vs.\ Rural Students (Tier A)]{sky}
An education researcher randomly samples students in urban and rural schools to compare
mean daily screen time (hours). Results:
Urban: $n_1 = 28$, $\bar x_1 = 5.4$ hr, $s_1 = 1.6$ hr. \quad
Rural: $n_2 = 25$, $\bar x_2 = 4.1$ hr, $s_2 = 1.3$ hr.
Both sample distributions are roughly symmetric with no outliers.

\begin{enumerate}[label=\alph*.]
  \item \textbf{Verify both conditions.}
        \begin{itemize}
          \item Random / Independence: \writeline
          \item Normal: \writeline
        \end{itemize}

  \item \textbf{Compute} $SE$ and construct a 95\% confidence interval for $\mu_{\text{urban}} - \mu_{\text{rural}}$.
        Technology gives $df \approx 50$, $t^* \approx 2.009$.

        $SE = \sqrt{\dfrac{(\blank{1.5cm})^2}{\blank{1.5cm}} + \dfrac{(\blank{1.5cm})^2}{\blank{1.5cm}}} \approx$ \blank{2.5cm}

        $CI = ($ \blank{2cm} $\pm$ \blank{2cm} $\times$ \blank{2cm} $) = ($ \blank{3cm} , \blank{3cm} $)$

  \item \textbf{Interpret} the interval in context and state what it implies about urban vs.\ rural screen time.
        \writelines{3}
\end{enumerate}
\end{scenariobox}

\begin{scenariobox}[Scenario 3 --- Compare Two Intervals (Tier E)]{sky}
A researcher computes two separate 95\% confidence intervals:
\begin{itemize}
  \item Group comparison A: $\mu_1 - \mu_2$: $(-0.3, \; 4.1)$ minutes
  \item Group comparison B: $\mu_3 - \mu_4$: $(2.8, \; 9.4)$ minutes
\end{itemize}

\begin{enumerate}[label=\alph*.]
  \item For interval A, does the data provide convincing evidence that $\mu_1 \ne \mu_2$?
        Explain. \writelines{2}

  \item For interval B, does the data provide convincing evidence that $\mu_3 > \mu_4$?
        Explain. \writelines{2}

  \item Interval B is wider than interval A. Give \textbf{two} possible reasons why one
        interval might be wider than another. \writelines{3}
\end{enumerate}
\end{scenariobox}

\end{document}
